% ============================================================
%  C: Como Programar -- 6ª Edição | Deitel & Deitel
%  Capítulo 7 -- Ponteiros em C
%  EXERCÍCIOS DE AUTORREVISÃO (7.1 -- 7.6)
%  Abordagem incremental: Caps. 1 + 2 + 3 + 4 + 5 + 6 + 7
% ============================================================
\documentclass[12pt,a4paper]{article}
\usepackage[utf8]{inputenc}
\usepackage[T1]{fontenc}
\usepackage[brazil]{babel}
\usepackage{geometry}
\geometry{top=2.5cm,bottom=2.5cm,left=3cm,right=3cm}
\usepackage{parskip}\setlength{\parskip}{6pt}\setlength{\parindent}{0pt}
\usepackage{xcolor,tcolorbox,enumitem,listings,fancyhdr,titlesec,hyperref,booktabs,array}
\tcbuselibrary{skins,breakable}

\definecolor{deitelblue}{RGB}{0,70,127}
\definecolor{deitelgray}{RGB}{240,242,245}
\definecolor{deitelgreen}{RGB}{0,110,60}
\definecolor{deitellightblue}{RGB}{220,235,250}
\definecolor{deitelred}{RGB}{160,20,20}
\definecolor{deitellorange}{RGB}{200,90,0}

\newtcolorbox{boaprática}[1][]{colback=deitellightblue,colframe=deitelblue,
  fonttitle=\bfseries\small,title={Boa prática de programação},breakable,#1}
\newtcolorbox{erroco}[1][]{colback=red!5,colframe=deitelred,
  fonttitle=\bfseries\small,title={Erro comum de programação},breakable,#1}
\newtcolorbox{prev}[1][]{colback=yellow!8,colframe=deitellorange,
  fonttitle=\bfseries\small,title={Dica de prevenção de erro},breakable,#1}
\newtcolorbox{respostabox}[1][]{colback=deitelgray,colframe=deitelblue!60,
  fonttitle=\bfseries\small\color{deitelblue},title={Resposta detalhada},breakable,#1}
\newtcolorbox{enunciadoBox}[1][]{colback=white,colframe=deitelblue,
  fonttitle=\bfseries,title={#1},breakable}
\newtcolorbox{saida}[1][]{colback=black!88,colframe=deitelblue,
  fontupper=\ttfamily\small\color{green!70},title={\small\color{white}Saída do programa},breakable,#1}

\setlist[enumerate,1]{label=\textbf{\alph*)},leftmargin=2em}
\setlist[itemize]{leftmargin=2em}

\lstset{
  basicstyle=\ttfamily\small,backgroundcolor=\color{deitelgray},
  frame=single,framerule=0.4pt,rulecolor=\color{deitelblue!40},
  breaklines=true,showstringspaces=false,
  keywordstyle=\color{deitelblue}\bfseries,
  commentstyle=\color{deitelgreen}\itshape,
  stringstyle=\color{deitelred},
  language=C,numbers=left,numberstyle=\tiny\color{gray},
  xleftmargin=18pt,xrightmargin=5pt,stepnumber=1,
}

\pagestyle{fancy}\fancyhf{}
\fancyhead[L]{\small\color{deitelblue}\textbf{C: Como Programar -- 6ª ed. | Deitel \& Deitel}}
\fancyhead[R]{\small\color{deitelblue}Cap. 7 -- Autorrevisão}
\fancyfoot[C]{\small\thepage}
\renewcommand{\headrulewidth}{0.4pt}

\titleformat{\section}[block]{\large\bfseries\color{deitelblue}}{\thesection.}{0.5em}{}[\titlerule]
\titleformat{\subsection}[block]{\normalsize\bfseries\color{deitelblue!80}}{}{}{}

\hypersetup{colorlinks=true,linkcolor=deitelblue}

\begin{document}

\begin{titlepage}
  \centering\vspace*{2cm}
  {\color{deitelblue}\rule{\linewidth}{2pt}}\\[0.4cm]
  {\huge\bfseries\color{deitelblue}C: Como Programar}\\[0.2cm]
  {\large\color{deitelblue}6ª Edição $\cdot$ Paul Deitel \& Harvey Deitel}\\[0.2cm]
  {\color{deitelblue}\rule{\linewidth}{2pt}}\\[1cm]
  {\LARGE\bfseries Capítulo 7}\\[0.3cm]
  {\Large\color{deitelblue!80}Ponteiros em C}\\[0.8cm]
  \begin{tcolorbox}[colback=deitellightblue,colframe=deitelblue,width=0.85\linewidth]
    \centering{\large\bfseries Exercícios de Autorrevisão (7.1--7.6)}\\[0.2cm]
    {\normalsize Resolvidos passo a passo, abordagem incremental\\
    Capítulos 1 + 2 + 3 + 4 + 5 + 6 + 7}
  \end{tcolorbox}
\end{titlepage}

\section*{Construindo sobre os Capítulos Anteriores}

No Capítulo~7, abordamos um dos tópicos mais poderosos -- e desafiadores -- de C:
os \textbf{ponteiros}. Construindo sobre os fundamentos dos Capítulos 1--6, incorporamos:

\begin{itemize}[noitemsep]
  \item \textbf{Ponteiros} -- variáveis que armazenam endereços de memória
  \item \textbf{Operador de endereço \texttt{\&}} -- obtém o endereço de uma variável
  \item \textbf{Operador de desreferência \texttt{*}} -- acessa o valor apontado
  \item \textbf{Chamada por referência simulada} -- passar ponteiros para funções
  \item \textbf{Aritmética de ponteiros} -- somar/subtrair inteiros a ponteiros
  \item \textbf{Equivalência entre arrays e ponteiros} -- nomes de arrays como ponteiros constantes
  \item \textbf{Ponteiros e \texttt{const}} -- 4 combinações de privilégios
  \item \textbf{Ponteiros para funções} -- passar e armazenar comportamento
  \item \textbf{Strings como arrays de \texttt{char}} -- preparação para o Cap.~8
\end{itemize}

\begin{boaprática}
  Ponteiros são a chave para entender a memória em C. Domine os operadores
  \texttt{\&} (endereço) e \texttt{*} (desreferência), a equivalência array/ponteiro
  e a aritmética de ponteiros. Esses três tópicos resolvem 90\% das dificuldades
  com ponteiros.
\end{boaprática}

\tableofcontents\newpage

% ============================================================
\section{Exercício 7.1 -- Preenchimento de Lacunas}
% ============================================================

\subsection*{Item (a)}
\begin{respostabox}
\textbf{Resposta:} Uma variável de ponteiro contém como valor o(a)
\textcolor{deitelblue}{\textbf{endereço}} de outra variável.

Um ponteiro \textbf{não} contém um valor de dados (como 42 ou 3.14); ele contém
o \textit{endereço de memória} onde um valor está armazenado. Em C, declaramos
ponteiros com asterisco antes do nome:
\begin{lstlisting}
int *ptr;          /* ptr e um ponteiro para int */
int x = 42;
ptr = &x;          /* ptr agora contem o ENDERECO de x */
printf( "%d\n", *ptr );  /* *ptr desreferencia: imprime 42 */
\end{lstlisting}

\textbf{Analogia:} se uma variável é uma ``caixa'' com um valor dentro, um ponteiro
é um ``papel com o número da caixa'' -- não a caixa em si, mas o caminho para
encontrá-la.
\end{respostabox}

\subsection*{Item (b)}
\begin{respostabox}
\textbf{Resposta:} Os três valores para inicializar um ponteiro são
\textcolor{deitelblue}{\textbf{0}}, \textcolor{deitelblue}{\textbf{NULL}} e
\textcolor{deitelblue}{\textbf{um endereço}}.

\begin{itemize}[noitemsep]
  \item \textbf{0} -- valor inteiro literal zero; tratado como ponteiro nulo em contexto de ponteiro
  \item \textbf{NULL} -- constante simbólica definida em \texttt{<stddef.h>} e outros
  cabeçalhos; representa ``nenhum endereço''
  \item \textbf{Um endereço} -- de outra variável, obtido com \texttt{\&}
\end{itemize}

\begin{boaprática}
  Use \texttt{NULL} (não 0) ao inicializar ponteiros, pois deixa a intenção
  explícita: \texttt{int *ptr = NULL;}. Inicializar ponteiros é fundamental --
  ponteiros não inicializados (``ponteiros perdidos'') causam comportamento
  imprevisível.
\end{boaprática}
\end{respostabox}

\subsection*{Item (c)}
\begin{respostabox}
\textbf{Resposta:} O único inteiro que pode ser atribuído a um ponteiro é
\textcolor{deitelblue}{\textbf{0}}.

\begin{lstlisting}
int *ptr;
ptr = 0;       /* OK: equivale a NULL */
ptr = NULL;    /* OK: forma preferida  */
ptr = 100;     /* ERRO (em compiladores modernos) */
\end{lstlisting}

A regra existe para evitar atribuições acidentais de números ``mágicos'' a ponteiros
-- algo que quase certamente causaria falha de segmentação. O zero é a exceção por
ter um significado especial como ``endereço inválido''.
\end{respostabox}

\newpage

% ============================================================
\section{Exercício 7.2 -- Verdadeiro ou Falso}
% ============================================================

\subsection*{Item (a) -- \texttt{\&} aplica-se a constantes, expressões e registradores}

\begin{tcolorbox}[colback=red!5,colframe=deitelred,title={\bfseries FALSO}]
O operador de endereço \texttt{\&} aplica-se \textbf{somente a variáveis}.
Não pode ser aplicado a:
\begin{itemize}[noitemsep]
  \item \textbf{Constantes:} \texttt{\&5} é inválido (constantes literais não têm endereço fixo).
  \item \textbf{Expressões:} \texttt{\&(x + y)} é inválido (expressões temporárias).
  \item \textbf{Variáveis \texttt{register}:} elas podem estar em registradores da
  CPU, que não possuem endereço na RAM.
\end{itemize}
\end{tcolorbox}

\subsection*{Item (b) -- Ponteiro \texttt{void} pode ser desreferenciado}

\begin{tcolorbox}[colback=red!5,colframe=deitelred,title={\bfseries FALSO}]
Um ponteiro \texttt{void *} \textbf{não} pode ser desreferenciado diretamente,
porque o compilador não sabe \textit{quantos bytes} ler nem \textit{como interpretá-los}.

Para desreferenciar, primeiro deve-se converter (\textit{cast}) o ponteiro para
o tipo apropriado:
\begin{lstlisting}
void *vp;
int x = 42;
vp = &x;

/* *vp; ERRO -- nao se pode desreferenciar void* */
int valor = *( ( int * )vp );  /* OK: converte para int* primeiro */
\end{lstlisting}

\textit{Por que existe \texttt{void *}?} É o ``ponteiro genérico'': pode receber
endereço de qualquer tipo, útil em funções de uso geral (como \texttt{malloc}, que
estudaremos no Cap.~12).
\end{tcolorbox}

\subsection*{Item (c) -- Ponteiros de tipos diferentes podem ser atribuídos sem cast}

\begin{tcolorbox}[colback=red!5,colframe=deitelred,title={\bfseries FALSO}]
Ponteiros de tipos diferentes \textbf{não} podem ser atribuídos uns aos outros
sem cast (com a única exceção do \texttt{void *}). Por exemplo:

\begin{lstlisting}
int   *ip;
float *fp;
ip = fp;            /* ERRO ou aviso de compilacao */
ip = ( int * )fp;   /* OK com cast (mas perigoso!) */
\end{lstlisting}

A regra existe porque a aritmética de ponteiros e a desreferenciação dependem do tipo
(\texttt{int *} avança 4 bytes; \texttt{float *} também 4 bytes; mas \texttt{double *}
avança 8 bytes, etc.). Atribuir entre tipos incompatíveis pode levar a interpretações
incorretas dos dados.

\textbf{Exceção:} \texttt{void *} pode ser atribuído a/de qualquer outro tipo de
ponteiro sem cast explícito.
\end{tcolorbox}

\newpage

% ============================================================
\section{Exercício 7.3 -- Array \texttt{numbers}}
% ============================================================

\begin{respostabox}
\begin{enumerate}[label=\textbf{\alph*)}]

\item \textbf{Definir e inicializar:}
\begin{lstlisting}
#define SIZE 10
float numbers[ SIZE ] = { 0.0, 1.1, 2.2, 3.3, 4.4,
                          5.5, 6.6, 7.7, 8.8, 9.9 };
\end{lstlisting}

\item \textbf{Ponteiro para \texttt{float}:}
\begin{lstlisting}
float *nPtr;
\end{lstlisting}

\item \textbf{Imprimir com notação de subscrito:}
\begin{lstlisting}
for ( i = 0; i < SIZE; ++i ) {
    printf( "%.1f ", numbers[ i ] );
}
\end{lstlisting}

\item \textbf{Atribuir endereço inicial a \texttt{nPtr} (duas formas):}
\begin{lstlisting}
nPtr = numbers;          /* nome do array e o endereco do 1o elemento */
nPtr = &numbers[ 0 ];    /* equivalente, mas explicito */
\end{lstlisting}

\item \textbf{Imprimir usando \texttt{nPtr} com notação ponteiro/deslocamento:}
\begin{lstlisting}
for ( i = 0; i < SIZE; ++i ) {
    printf( "%.1f ", *( nPtr + i ) );
}
\end{lstlisting}
\textit{Note:} \texttt{*(nPtr + i)} é equivalente a \texttt{nPtr[i]} ou \texttt{numbers[i]}.

\item \textbf{Imprimir usando o nome do array como ponteiro:}
\begin{lstlisting}
for ( i = 0; i < SIZE; ++i ) {
    printf( "%.1f ", *( numbers + i ) );
}
\end{lstlisting}

\item \textbf{Imprimir subscritando \texttt{nPtr}:}
\begin{lstlisting}
for ( i = 0; i < SIZE; ++i ) {
    printf( "%.1f ", nPtr[ i ] );
}
\end{lstlisting}

\item \textbf{Referenciar o elemento 4 de \textbf{quatro formas:}}
\begin{lstlisting}
numbers[ 4 ]        /* notacao de subscrito de array     */
*( numbers + 4 )    /* ponteiro/deslocamento com nome do array */
nPtr[ 4 ]           /* notacao de subscrito de ponteiro  */
*( nPtr + 4 )       /* ponteiro/deslocamento com nPtr    */
\end{lstlisting}

\item \textbf{Endereço de \texttt{nPtr + 8}:}\\
Endereço inicial: 1002500. Cada \texttt{float} ocupa 4 bytes. Aritmética:
\[
1002500 + 8 \times 4 = 1002500 + 32 = \mathbf{1002532}
\]
Valor armazenado em \texttt{numbers[8]}: \textbf{8.8}.

\item \textbf{Endereço após \texttt{nPtr -= 4}, supondo \texttt{nPtr} aponta para \texttt{numbers[5]}:}
\begin{itemize}[noitemsep]
  \item Endereço de \texttt{numbers[5]}: $1002500 + 5 \times 4 = 1002520$
  \item Após \texttt{nPtr -= 4}: $1002520 - 4 \times 4 = \mathbf{1002504}$
\end{itemize}
Esse endereço é \texttt{numbers[1]}, cujo valor é \textbf{1.1}.

\end{enumerate}

\begin{boaprática}
  As \textbf{quatro notações equivalentes} -- \texttt{numbers[i]}, \texttt{nPtr[i]},
  \texttt{*(numbers + i)}, \texttt{*(nPtr + i)} -- são uma das ideias mais
  poderosas de C. Domine-as e ponteiros deixam de parecer ``mágicos''.
\end{boaprática}
\end{respostabox}

\newpage

% ============================================================
\section{Exercício 7.4 -- Variáveis \texttt{number1}, \texttt{number2} e \texttt{fPtr}}
% ============================================================

\begin{respostabox}
Considere \texttt{float number1 = 7.3;} e \texttt{float number2;}.

\begin{enumerate}[label=\textbf{\alph*)}]

\item \textbf{Definir ponteiro \texttt{fPtr}:}
\begin{lstlisting}
float *fPtr;
\end{lstlisting}

\item \textbf{Atribuir endereço de \texttt{number1}:}
\begin{lstlisting}
fPtr = &number1;
\end{lstlisting}

\item \textbf{Imprimir valor apontado:}
\begin{lstlisting}
printf( "O valor de *fPtr e %f\n", *fPtr );
\end{lstlisting}

\item \textbf{Atribuir o valor apontado a \texttt{number2}:}
\begin{lstlisting}
number2 = *fPtr;
\end{lstlisting}

\item \textbf{Imprimir \texttt{number2}:}
\begin{lstlisting}
printf( "O valor de number2 e %f\n", number2 );
\end{lstlisting}

\item \textbf{Imprimir endereço de \texttt{number1} com \texttt{\%p}:}
\begin{lstlisting}
printf( "O endereco de number1 e %p\n", &number1 );
\end{lstlisting}

\item \textbf{Imprimir endereço armazenado em \texttt{fPtr}:}
\begin{lstlisting}
printf( "O endereco armazenado em fPtr e %p\n", fPtr );
\end{lstlisting}

\textbf{O valor impresso é o mesmo que \texttt{\&number1}?} \textbf{Sim!} Como
\texttt{fPtr = \&number1} no item (b), \texttt{fPtr} contém exatamente o endereço
de \texttt{number1}. Os dois \texttt{printf} produzirão a mesma saída.

\end{enumerate}

\textit{Nota sobre \texttt{\%p}:} é o especificador padrão para imprimir endereços
em formato hexadecimal. A representação exata depende do compilador (geralmente
\texttt{0x7fffe1a23bdc} ou similar).
\end{respostabox}

\newpage

% ============================================================
\section{Exercício 7.5 -- Protótipos e Cabeçalhos}
% ============================================================

\begin{respostabox}
\begin{enumerate}[label=\textbf{\alph*)}]

\item \textbf{Cabeçalho de \texttt{exchange}}: dois ponteiros para \texttt{float}, retorna \texttt{void}.
\begin{lstlisting}
void exchange( float *x, float *y )
\end{lstlisting}

\item \textbf{Protótipo de \texttt{exchange}}: adicione \texttt{;}.
\begin{lstlisting}
void exchange( float *x, float *y );
\end{lstlisting}

\item \textbf{Cabeçalho de \texttt{evaluate}}: recebe \texttt{int x} e um ponteiro
para a função \texttt{poly} (que recebe \texttt{int} e retorna \texttt{int}); retorna \texttt{int}.
\begin{lstlisting}
int evaluate( int x, int (*poly)( int ) )
\end{lstlisting}

\item \textbf{Protótipo de \texttt{evaluate}}:
\begin{lstlisting}
int evaluate( int x, int (*poly)( int ) );
\end{lstlisting}

\end{enumerate}

\textbf{Sintaxe importante:} a expressão \texttt{int (*poly)(int)} é lida como
``\texttt{poly} é um ponteiro para função que recebe um \texttt{int} e retorna
um \texttt{int}''. Os parênteses ao redor de \texttt{*poly} são \textbf{essenciais}:
sem eles, \texttt{int *poly(int)} seria interpretado como ``função que retorna
\texttt{int *}'' -- algo totalmente diferente!

\begin{boaprática}
  Ponteiros para funções são poderosos porque permitem \textit{passar comportamento}
  como argumento -- a base de \textit{callbacks}, \textit{strategy pattern} e
  programação funcional. Já vimos um exemplo prático: \texttt{qsort} aceita um
  ponteiro para função de comparação.
\end{boaprática}
\end{respostabox}

\newpage

% ============================================================
\section{Exercício 7.6 -- Identificar e Corrigir Erros}
% ============================================================

\textbf{Considere as declarações:}
\begin{lstlisting}
int *zPtr;                            /* nao inicializado! */
int *aPtr = NULL;
void *sPtr = NULL;
int number, i;
int z[ 5 ] = { 1, 2, 3, 4, 5 };
sPtr = z;
\end{lstlisting}

\subsection*{Item (a)}
\begin{respostabox}
\textbf{Código com erro:} \texttt{++zptr;}

\textbf{Erros:}
\begin{enumerate}[noitemsep]
  \item \texttt{zptr} (minúsculo) é diferente de \texttt{zPtr}; C é \textit{case-sensitive}.
  \item Mesmo escrevendo \texttt{++zPtr}, há erro: \texttt{zPtr} nunca foi inicializado.
  Incrementar um ponteiro indefinido leva a comportamento indefinido.
\end{enumerate}

\textbf{Correção:}
\begin{lstlisting}
zPtr = z;     /* inicializa antes de usar */
++zPtr;       /* agora aponta para z[1] */
\end{lstlisting}

\begin{erroco}
  Usar um ponteiro \textbf{não inicializado} é uma das principais causas de
  falhas em C. Sempre inicialize ponteiros (com \texttt{NULL}, \texttt{0} ou
  um endereço válido) antes de usá-los.
\end{erroco}
\end{respostabox}

\subsection*{Item (b)}
\begin{respostabox}
\textbf{Código com erro:}
\begin{lstlisting}
/* usa o ponteiro para obter o primeiro valor do array */
number = zPtr;
\end{lstlisting}

\textbf{Erro:} \texttt{zPtr} é um \textit{endereço} (ponteiro), não um valor inteiro.
Atribuí-lo a \texttt{number} (que é \texttt{int}) gera aviso ou erro.

\textbf{Correção:} desreferenciar com \texttt{*}:
\begin{lstlisting}
number = *zPtr;  /* obtem o valor apontado */
\end{lstlisting}
\end{respostabox}

\subsection*{Item (c)}
\begin{respostabox}
\textbf{Código com erro:}
\begin{lstlisting}
number = *zPtr[ 2 ];
\end{lstlisting}

\textbf{Erro:} \texttt{zPtr[2]} já é um \textit{valor} \texttt{int} (não um ponteiro).
Aplicar \texttt{*} a um inteiro é incorreto.

\textbf{Correção:}
\begin{lstlisting}
number = zPtr[ 2 ];  /* equivale a *(zPtr + 2) */
\end{lstlisting}
\end{respostabox}

\subsection*{Item (d)}
\begin{respostabox}
\textbf{Código com erro:}
\begin{lstlisting}
for ( i = 0; i <= 5; i++ ) {
    printf( "%d ", zPtr[ i ] );
}
\end{lstlisting}

\textbf{Erro:} \texttt{i <= 5} acessa \texttt{zPtr[5]}, que está \textbf{fora dos
limites} do array \texttt{z} (índices válidos: 0 a 4).

\textbf{Correção:}
\begin{lstlisting}
for ( i = 0; i < 5; ++i ) {  /* < em vez de <= */
    printf( "%d ", zPtr[ i ] );
}
\end{lstlisting}
\end{respostabox}

\subsection*{Item (e)}
\begin{respostabox}
\textbf{Código com erro:} \texttt{number = *sPtr;}

\textbf{Erro:} \texttt{sPtr} é \texttt{void *} -- não pode ser desreferenciado
diretamente sem cast.

\textbf{Correção:}
\begin{lstlisting}
number = *( ( int * )sPtr );  /* cast para int* antes */
\end{lstlisting}
\end{respostabox}

\subsection*{Item (f)}
\begin{respostabox}
\textbf{Código com erro:} \texttt{++z;}

\textbf{Erro:} O nome de um array (\texttt{z}) é um \textit{ponteiro constante} --
seu valor não pode ser modificado.

\textbf{Correção:} usar uma variável de ponteiro auxiliar:
\begin{lstlisting}
int *p = z;
++p;            /* agora p aponta para z[1]; z permanece valido */
\end{lstlisting}

Alternativamente, usar subscrito direto: \texttt{z[1]} para acessar o segundo elemento.
\end{respostabox}

\newpage

% ============================================================
\section*{Gabarito Rápido -- Capítulo 7: Autorrevisão}

\begin{tcolorbox}[colback=deitellightblue,colframe=deitelblue,title={\bfseries\large Respostas Resumidas}]

\textbf{7.1:}
\begin{itemize}[noitemsep]
  \item[a)] endereço \quad b) 0, NULL, um endereço \quad c) 0
\end{itemize}

\medskip\textbf{7.2 -- Verdadeiro/Falso (todos falsos):}
\begin{itemize}[noitemsep]
  \item[a)] \textcolor{deitelred}{\textbf{Falso}} -- \texttt{\&} só em variáveis (não \texttt{register})
  \item[b)] \textcolor{deitelred}{\textbf{Falso}} -- \texttt{void *} não pode ser desreferenciado
  \item[c)] \textcolor{deitelred}{\textbf{Falso}} -- precisa de cast (exceto \texttt{void *})
\end{itemize}

\medskip\textbf{7.3:} arrays e ponteiros; quatro notações equivalentes;
endereços calculados pela aritmética de ponteiros (deslocamento $\times$ tamanho do tipo).

\medskip\textbf{7.4:} \texttt{fPtr = \&number1}, desreferência com \texttt{*fPtr},
\texttt{\%p} para endereços (idênticos para \texttt{\&number1} e \texttt{fPtr}).

\medskip\textbf{7.5:} sintaxe de ponteiros para funções: \texttt{int (*poly)(int)}.

\medskip\textbf{7.6:} 6 erros corrigidos: ponteiro não inicializado, desreferência
faltando, desreferência incorreta, \texttt{<=} em vez de \texttt{<}, falta cast em
\texttt{void *}, e modificação de nome de array.

\end{tcolorbox}

\end{document}
